\documentclass[10pt,a4paper]{amsart}
\usepackage{amsmath,amssymb,amsthm,graphicx}
% Retained preamble for the Stacks Project chapter snapshot in this directory.
% The chapter's own preamble.tex is a copy of the shared Stacks Project preamble; the block
% below is sliced MECHANICALLY, line for line, from it. Nothing was added, removed or
% redefined; the wrapper class, the xr-hyper/hyperref wiring and the externaldocument list
% are omitted because the assembler supplies the document class.
%
% lines taken from preamble.tex: 3, 5, 9, 11, 12, 156, 157, 158, ... (51 lines total)
%
% package coverage restored from the chapter's own preamble (2026-10-13): xy, xy-2cell

\usepackage[all]{xy}
\xyoption{2cell}
\UseAllTwocells

\newenvironment{reference}{\comment}{\endcomment}
\newenvironment{slogan}{\comment}{\endcomment}
\newenvironment{history}{\comment}{\endcomment}
\theoremstyle{plain}
\newtheorem{theorem}[subsection]{Theorem}
\newtheorem{proposition}[subsection]{Proposition}
\newtheorem{lemma}[subsection]{Lemma}
\theoremstyle{definition}
\newtheorem{definition}[subsection]{Definition}
\newtheorem{example}[subsection]{Example}
\newtheorem{exercise}[subsection]{Exercise}
\newtheorem{situation}[subsection]{Situation}
\theoremstyle{remark}
\newtheorem{remark}[subsection]{Remark}
\newtheorem{remarks}[subsection]{Remarks}
\numberwithin{equation}{subsection}
\def\lim{\mathop{\mathrm{lim}}\nolimits}
\def\colim{\mathop{\mathrm{colim}}\nolimits}
\def\Spec{\mathop{\mathrm{Spec}}}
\def\Hom{\mathop{\mathrm{Hom}}\nolimits}
\def\Ext{\mathop{\mathrm{Ext}}\nolimits}
\def\SheafHom{\mathop{\mathcal{H}\!\mathit{om}}\nolimits}
\def\SheafExt{\mathop{\mathcal{E}\!\mathit{xt}}\nolimits}
\def\Sch{\mathit{Sch}}
\def\Mor{\mathop{\mathrm{Mor}}\nolimits}
\def\Ob{\mathop{\mathrm{Ob}}\nolimits}
\def\Sh{\mathop{\mathit{Sh}}\nolimits}
\def\NL{\mathop{N\!L}\nolimits}
\def\CH{\mathop{\mathrm{CH}}\nolimits}
\def\proetale{{pro\text{-}\acute{e}tale}}
\def\etale{{\acute{e}tale}}
\def\QCoh{\mathit{QCoh}}
\def\Ker{\mathop{\mathrm{Ker}}}
\def\Im{\mathop{\mathrm{Im}}}
\def\Coker{\mathop{\mathrm{Coker}}}
\def\Coim{\mathop{\mathrm{Coim}}}
\DeclareMathSymbol{\boxtimes}{\mathbin}{AMSa}{"02}
\def\QCohstack{\mathcal{QC}\!\mathit{oh}}
\def\Cohstack{\mathcal{C}\!\mathit{oh}}
\def\Spacesstack{\mathcal{S}\!\mathit{paces}}
\def\Quotfunctor{\mathrm{Quot}}
\def\Hilbfunctor{\mathrm{Hilb}}
\def\Curvesstack{\mathcal{C}\!\mathit{urves}}
\def\Polarizedstack{\mathcal{P}\!\mathit{olarized}}
\def\Complexesstack{\mathcal{C}\!\mathit{omplexes}}
\def\Pic{\mathop{\mathrm{Pic}}\nolimits}
\def\Picardstack{\mathcal{P}\!\mathit{ic}}
\def\Picardfunctor{\mathrm{Pic}}
\def\Deformationcategory{\mathcal{D}\!\mathit{ef}}
\usepackage{cleveref}
\crefname{originallemma}{original lemma}{original lemmas}
\Crefname{originallemma}{Original lemma}{Original lemmas}
\makeatletter
\newlabel{lemma-presentation-quasi-compact}{{0265}{source}{}{sourcepointer-1}{}}
\newlabel{lemma-presentation-quasi-compact@cref}{{[originallemma][0][]0265}{[1][0][]source}{}{}{}}
\makeatother
\crefname{originallemma}{original lemma}{original lemmas}
\Crefname{originallemma}{Original lemma}{Original lemmas}
\makeatletter
\newlabel{lemma-representable-morphisms-spaces-property}{{02WZ}{source}{}{sourcepointer-2}{}}
\newlabel{lemma-representable-morphisms-spaces-property@cref}{{[originallemma][0][]02WZ}{[1][0][]source}{}{}{}}
\makeatother
\crefname{originallemma}{original lemma}{original lemmas}
\Crefname{originallemma}{Original lemma}{Original lemmas}
\makeatletter
\newlabel{morphisms-lemma-finite-fibre}{{02NG}{source}{}{sourcepointer-3}{}}
\newlabel{morphisms-lemma-finite-fibre@cref}{{[originallemma][0][]02NG}{[1][0][]source}{}{}{}}
\makeatother
\crefname{originallemma}{original lemma}{original lemmas}
\Crefname{originallemma}{Original lemma}{Original lemmas}
\makeatletter
\newlabel{morphisms-lemma-permanence-finite-type}{{01T8}{source}{}{sourcepointer-4}{}}
\newlabel{morphisms-lemma-permanence-finite-type@cref}{{[originallemma][0][]01T8}{[1][0][]source}{}{}{}}
\makeatother
\crefname{originalremark}{original remark}{original remarks}
\Crefname{originalremark}{Original remark}{Original remarks}
\makeatletter
\newlabel{remark-list-properties-fpqc-local-base}{{02WH}{source}{}{sourcepointer-5}{}}
\newlabel{remark-list-properties-fpqc-local-base@cref}{{[originalremark][0][]02WH}{[1][0][]source}{}{}{}}
\makeatother
\crefname{originalremark}{original remark}{original remarks}
\Crefname{originalremark}{Original remark}{Original remarks}
\makeatletter
\newlabel{remark-list-properties-stable-base-change}{{02WF}{source}{}{sourcepointer-6}{}}
\newlabel{remark-list-properties-stable-base-change@cref}{{[originalremark][0][]02WF}{[1][0][]source}{}{}{}}
\makeatother
\crefname{originallemma}{original lemma}{original lemmas}
\Crefname{originallemma}{Original lemma}{Original lemmas}
\makeatletter
\newlabel{schemes-lemma-monomorphism-separated}{{01L4}{source}{}{sourcepointer-7}{}}
\newlabel{schemes-lemma-monomorphism-separated@cref}{{[originallemma][0][]01L4}{[1][0][]source}{}{}{}}
\makeatother
\title{The diagonal of an algebraic space is locally of finite type, a monomorphism, separated and locally quasi-finite}
\author{The Stacks Project Authors}
\date{stacks-project:spaces.tex:local-snapshot-20261012; GNU FDL 1.2 or later}
\begin{document}
\maketitle
\setcounter{subsection}{1}
\noindent\textit{Editorial note.} All delivered mathematical text is the author's own, copied from the pinned Stacks Project chapter file spaces.tex; no statement or proof step was rewritten, reordered or removed.\par\medskip
\noindent\textit{Editorial note.} The proof contains one xymatrix display; it is part of the counted proof span and is preserved verbatim, with the chapter's own Xy-pic setup retained in the preamble.\par\medskip
\noindent\textit{Editorial note.} The seven references of the counted spans are recorded as non-load-bearing navigation pointers; the argument that consumes them stays inside the counted proof span.\par\medskip
\noindent\textit{Editorial note.} Preamble: sources/stacks-spaces/retained-preamble.tex, sliced mechanically from the chapter's own preamble.tex.\par\medskip
\noindent\textit{Allowed standard prerequisites (editorial).}\begin{itemize}
\item The cartesian square computing the diagonal from a presentation. used at proof lines 2311-2318. Reduces the properties of Delta to those of j.
\item Properties of representable morphisms of spaces. Lemma lemma-representable-morphisms-spaces-property (used at proof lines 2319-2321), together with Remarks remark-list-properties-stable-base-change and remark-list-properties-fpqc-local-base. Transfers the four properties between Delta and j.
\item A monomorphism of schemes is separated. Schemes, Lemma schemes-lemma-monomorphism-separated, used at proof lines 2325-2327. Gives separatedness of j.
\item Permanence of finite type and finite fibres. Morphisms, Lemmas morphisms-lemma-permanence-finite-type and morphisms-lemma-finite-fibre, used at proof lines 2328-2334. Give local finiteness of type and local quasi-finiteness of j.
\end{itemize}\medskip
\ifdefined\hypertarget\hypertarget{sourcepointer-1}{}\fi
\noindent\textit{Source reference.} Original lemma 0265: focused-seeds/AGP/sources/stacks-snapshot-20261011/tags.txt, source lines 2838--2838. Navigation pointer into the Stacks Project; the standard fact it supplies is applied inside the counted proof span and the argument resting on it stays there.\par\medskip
\ifdefined\hypertarget\hypertarget{sourcepointer-2}{}\fi
\noindent\textit{Source reference.} Original lemma 02WZ: focused-seeds/AGP/sources/stacks-snapshot-20261011/tags.txt, source lines 3822--3822. Navigation pointer into the Stacks Project; the standard fact it supplies is applied inside the counted proof span and the argument resting on it stays there.\par\medskip
\ifdefined\hypertarget\hypertarget{sourcepointer-3}{}\fi
\noindent\textit{Source reference.} Original lemma 02NG: focused-seeds/AGP/sources/stacks-snapshot-20261011/tags.txt, source lines 3479--3479. Navigation pointer into the Stacks Project; the standard fact it supplies is applied inside the counted proof span and the argument resting on it stays there.\par\medskip
\ifdefined\hypertarget\hypertarget{sourcepointer-4}{}\fi
\noindent\textit{Source reference.} Original lemma 01T8: focused-seeds/AGP/sources/stacks-snapshot-20261011/tags.txt, source lines 2365--2365. Navigation pointer into the Stacks Project; the standard fact it supplies is applied inside the counted proof span and the argument resting on it stays there.\par\medskip
\ifdefined\hypertarget\hypertarget{sourcepointer-5}{}\fi
\noindent\textit{Source reference.} Original remark 02WH: focused-seeds/AGP/sources/stacks-snapshot-20261011/tags.txt, source lines 3804--3804. Navigation pointer into the Stacks Project; the standard fact it supplies is applied inside the counted proof span and the argument resting on it stays there.\par\medskip
\ifdefined\hypertarget\hypertarget{sourcepointer-6}{}\fi
\noindent\textit{Source reference.} Original remark 02WF: focused-seeds/AGP/sources/stacks-snapshot-20261011/tags.txt, source lines 3802--3802. Navigation pointer into the Stacks Project; the standard fact it supplies is applied inside the counted proof span and the argument resting on it stays there.\par\medskip
\ifdefined\hypertarget\hypertarget{sourcepointer-7}{}\fi
\noindent\textit{Source reference.} Original lemma 01L4: focused-seeds/AGP/sources/stacks-snapshot-20261011/tags.txt, source lines 2073--2073. Navigation pointer into the Stacks Project; the standard fact it supplies is applied inside the counted proof span and the argument resting on it stays there.\par\medskip
\begin{lemma}
\label{lemma-properties-diagonal}
Let $S$ be a scheme contained in $\Sch_{fppf}$.
Let $F$ be an algebraic space over $S$.
Let $\Delta : F \to F \times F$ be the diagonal morphism.
Then
\begin{enumerate}
\item $\Delta$ is locally of finite type,
\item $\Delta$ is a monomorphism,
\item $\Delta$ is separated, and
\item $\Delta$ is locally quasi-finite.
\end{enumerate}
\end{lemma}

\begin{proof}
Let $F = U/R$ be a presentation of $F$.
As in the proof of Lemma \ref{lemma-presentation-quasi-compact} the diagram
$$
\xymatrix{
R \ar[r] \ar[d]_j & F \ar[d]^\Delta \\
U \times_S U \ar[r] & F \times F
}
$$
is cartesian. Hence according to
Lemma \ref{lemma-representable-morphisms-spaces-property}
it suffices to show that $j$ has the properties listed in the lemma.
(Note that each of the properties (1) -- (4) occur in the lists
of Remarks \ref{remark-list-properties-stable-base-change}
and \ref{remark-list-properties-fpqc-local-base}.)
Since $j$ is an equivalence relation it is a monomorphism.
Hence it is separated by
Schemes, Lemma \ref{schemes-lemma-monomorphism-separated}.
As $R$ is an \'etale equivalence relation we see that
$s, t : R \to U$ are \'etale. Hence $s, t$ are locally of finite
type. Then it follows from
Morphisms, Lemma \ref{morphisms-lemma-permanence-finite-type} that
$j$ is locally of finite type. Finally, as it is a monomorphism
its fibres are finite. Thus we conclude that it is locally quasi-finite by
Morphisms, Lemma \ref{morphisms-lemma-finite-fibre}.
\end{proof}

\end{document}
